\subsection{The class diagram and the manuscript corrections}
\label{subsec:manuscript-audit}

\begin{proposition}[Known inclusions; two adjacent comparisons unresolved]
\label{D:class-diagram}
On nonempty connected small categories, the verified conclusions are
\[
 \{\text{filtered}\}\subsetneq\{\text{exact}\}
 \subseteq\{\text{total}\}\subseteq\{\text{quasi-filtered}\}.
\]
The exact and quasi-filtered classes differ. This record does not decide
either exact $=$ total or total $=$ quasi-filtered separately.
\end{proposition}
\begin{proof}
We use the checked implications R1--R4 in the supplied prompt. The
category $\Delta_{\mathrm{inj}}$ of nonempty finite linearly ordered
sets and increasing injections is exact by
\citet[Theorem~3]{IsbellMitchell1973}, and is not filtered: its two
maps from a singleton to a two-element chain cannot be equalized by
an injection. Exactness implies totality by R2. The quasi-filtered,
nonexact category of \citet[Construction~C.1 and Theorem~C.2]{AdamekKoubekReiterman1981}
separates the outer two classes. Part~C records the checked proof of
that theorem and the repairs used for its parameters. Its totality
has not been proved or refuted in this return. Thus its nonexactness
does not identify which adjacent inclusion is strict.

We also read the full interpolation verification in TOTASPH C.8--C.10.
The integral coordinate change $x_i=z_i-z_{i+1}$ reduces compatible
coface data to coordinate-hyperplane interpolation, whose alternating
sum preserves homogeneity and augmentation one. The prescriptions
$D_iC_{n,k}=C_{n-1,2k}\prod_{l\ne i}y_l^k$ agree on adjacent
intersections directly and on nonadjacent intersections by the
inductive normalized-face identity. Gap encoding turns the resulting
polynomial into an affine arrow equalizing all cofaces. The simplicial
identities then equalize every ordinary parallel family. This confirms
the exactness used here without treating the unopened 1974 body as a proof.

We re-read the needed proof steps of R1, R2 and R4 in TOTASPH
C.3--C.6 and C.11 and TOTEX A.1--A.4. In the sufficiency direction
of R1 used below, for an ordinary final $f:A\to B$ the comma
category $b\downarrow f$ is exactly the category of elements of
$a\mapsto B(b,f(a))$. Thus the action condition makes each comma
category weakly contractible, and Theorem~A applies to the specified
map $Bf$. For the implication exact $\Rightarrow$ total in R2,
the functor
\[
 \Sigma_H:\Ab^{\El_AH}\longrightarrow\Ab^A,
 \qquad (\Sigma_H F)(a)=\bigoplus_{h\in H(a)}F(a,h),
\]
is exact and has a canonical colimit comparison
$\colim_{\El_AH}F\cong\colim_A\Sigma_H F$. Hence all action
components have exact colimits, and their constant coefficient bar
complexes are integrally acyclic. The inherited-action construction
in the next proof lets us apply the same conclusion to every
connected covering of every such component. In particular the
cover belonging to a cyclic subgroup $K$ of its fundamental group
has $H_1=K_{\mathrm{ab}}=K=0$. Every cyclic subgroup is trivial;
the component is simply connected and acyclic, and Hurewicz gives
weak contractibility. This verifies the crucial integral-to-homotopy
step without presupposing the converse to R2.

For R4, if two objects have no common successor, their disjoint
principal cosieves $U$ and $V$ give the action formed by identifying
the two constant values precisely on $U\amalg V$. Its nerve is
two copies of $NA$ glued on $N(U\amalg V)$, because arrows cannot
leave a cosieve. Each cosieve is connected. Mayer--Vietoris gives
$H_1=\Z$, contradicting totality. Applying the same argument to
$p/A=\El_A A(p,-)$ completes every span out of $p$ to a square.
For $u,v:p\to q$ the quotient action
$H=A(p,-)/(cu\sim cv)$ has free presentation
$P_q\xrightarrow{v-u}P_p\to\Z[H]\to0$. If
$K=\ker(P_q\to P_p)$, the two short exact sequences obtained by
inserting its image give
$H_1(B\El_AH;\Z)=\Z/\aug(\colim_A K)$. Vanishing gives an
augmentation-one finite sum in $K$; common successors move its
finitely many supports to a single target. This is IV, together
with I proved above. In particular no step replaces a kernel in a
direct sum by an intersection of kernels.
\end{proof}

\begin{proposition}[Refuted conjecture; known R6 repairs of II ms.\ 36--36ter]
\label{D:manuscript-II}
The conjecture printed at II ms.\ 36ter,
\[
 X\text{ totalement }W_\infty\text{-asph\'erique}
 \Longrightarrow X\text{ filtrante},
\]
is false. For nonempty connected small categories its following
strengthenings are true:
\begin{enumerate}[label=(\roman*)]
\item a total sifted category is filtered;
\item for any fixed integer $n\geq1$, a total category $A$ is filtered
if $\hocolim_A K(M(a),n)$ is $n$-truncated for every covariant diagram
$M:A\to\Ab$.
\end{enumerate}
Filtered categories satisfy both additional conditions. These are the
checked R6 results, not a resolution of totality versus exactness.
\end{proposition}
\begin{proof}[Reading and proof check]
We read II ms.\ 36, 36bis and 36ter in the supplied
K\"unzer--Malgoire--Maltsiniotis transcription
\citep{GrothendieckDerivateurs}. The category
$\Delta_{\mathrm{inj}}$ just recalled refutes the conjecture. This
uses the published example and makes no claim of a new counterexample.
The original body of Isbell's 1974 paper was not obtained in this
session; its proof is therefore not represented as having been read.

There are two distinct assertions on ms.\ 36. The proposed
Proposition~2.19, which infers totality from finality of the diagonal,
is expressly footnoted ``Faux.'' The subsequent question asks whether
totality implies that final-diagonal condition, or at least PS~$1'$.
With the nonemptiness convention, finality of the diagonal is
siftedness. Totality does not imply siftedness, since (i) and
$\Delta_{\mathrm{inj}}$ would otherwise imply filteredness. It does
imply PS~$1'$: any two objects have a common successor. This follows
already from the span-completion condition I in R4. The printed
corollary claiming totality of $\Delta^{\mathrm{op}}$ is false as
well. For example, the simplicial circle, regarded as a functor
$\Delta^{\mathrm{op}}\to\Set$, has a category of elements whose
classifying space is weakly equivalent to the circle, by the usual
simplex-category realization comparison. Its first homology is
$\Z$, so it violates the action criterion R1.

For clarity we spell out the short proof of (i), also proved in the
supplied PRETOT Lemma~F.3 and rechecked in TOTEX D.4. Totality passes
to every connected component $C$ of $\El_A S$: for $T:C\to\Set$ the
functor
\[
 R(a)=\coprod_{\substack{s\in S(a)\\(a,s)\in C}}T(a,s)
 \quad(a\in A)
\]
has structure maps induced by those of $S$ and $T$, and there is an
actual isomorphism $\El_A R\cong\El_C T$. In particular such a
nonempty connected $C$ has common successors, by R4. If $A$ is
sifted, its cospan category
\[
 C_{p,q}=\El_A\bigl(A(p,-)\times A(q,-)\bigr)
\]
is nonempty and connected for every $p,q$. This is the criterion of
\citet[Example~15.2(c), proof in \S15.9]{GabrielUlmer1971}; both the
statement and that proof were opened in this session. Given
$u,v:p\to q$, a common successor in $C_{p,q}$ of
$(q,u,\id_q)$ and $(q,v,\id_q)$ supplies maps $f,g:q\to r$ with
$fu=gv$ and $f\id_q=g\id_q$. Hence $f=g$ and $fu=fv$.
Siftedness itself supplies common successors in $A$, so $A$ is
filtered. Every map in this argument is an ordinary arrow of $A$.

We also re-read the complete abelian and fixed-$n$ proofs of R6 in
the supplied TOTEX D.5 and D.7. Their relevant comparisons are as
follows. Write $G=\colim_A^{\Ab}M$. If $n\geq2$, totality gives
$BA\simeq *$; connectivity of the bar construction and its degree-$n$
edge map show that
\[
 X_M=\hocolim_A K(M(a),n)\longrightarrow K(G,n)
\]
is a weak equivalence under the assumed $n$-truncation. The reduced
bar spectral sequence has row $n$ equal to $M$, no lower rows and
zero row $n+1$, since $H_{n+1}(K(L,n);\Z)=0$ for every abelian
group $L$. Thus
\[
 L_1\colim_A M=H_{n+1}(X_M;\Z)=0.
\]
This proves exactness, rather than assuming the unresolved converse
to R2. Exactness then collapses the same homology spectral sequence
in every degree. Apply the resulting canonical homology comparison
to $M\oplus N$ and take the cross summand in degree $2n$ of
$K(M,n)\times K(N,n)$. The integral K\"unneth formula gives
\[
 \colim_A(M\otimes_\Z N)
 \xrightarrow{\ \cong\ }
 (\colim_A M)\otimes_\Z(\colim_A N).
\]
No Tor summand occurs in that cross degree, since both positive
indices are at least $n$. Taking $M=\Z[A(p,-)]$ and
$N=\Z[A(q,-)]$ makes this the free abelian map from
$\pi_0 C_{p,q}$ to a singleton. Consequently $C_{p,q}$ is nonempty
and connected. Hence $A$ is sifted, and (i) applies.

For $n=1$, common successors, already provided by totality, make
the colimit in groups of the diagram $M$ abelian. The fundamental
group presentation of $X_M=\hocolim_A BM$ therefore identifies
its fundamental group canonically with $G$; the transport part of
that presentation is killed by $\pi_1(BA)=0$. The assumed
$1$-truncation makes $X_M\to BG$ a weak equivalence. In the
homology spectral sequence, $E^\infty_{1,1}=L_1\colim_A M$:
the only potentially incoming differential is from
$H_3(BA)=0$. The edge map on second homology is the canonical
\[
 \colim_A\Lambda^2 M\longrightarrow\Lambda^2G.
\]
It is surjective because two representatives can be moved to a
common successor. Since $H_2(X_M)\to H_2(BG)$ is an isomorphism,
the edge filtration exhausts $H_2(X_M)$ and
$L_1\colim_A M=0$. Exactness follows. Now the edge map is an
isomorphism; the natural decomposition
$\Lambda^2(M\oplus N)=\Lambda^2M\oplus(M\otimes N)\oplus\Lambda^2N$
gives the same tensor comparison and siftedness. This checks the
argument credited in the supplied prompt to PRETOT F.6 as well as
TOTEX D.5. For the converse, filtered homotopy colimits preserve
$n$-truncation: the tests for homotopy groups and for equality of
their elements involve finite simplicial data.
\end{proof}

\begin{remark}[Known: the positive-truncation result retained from R6]
\label{D:positive-truncation-record}
For an arbitrary small category $A$, including the empty and
disconnected cases, put $\ell:A\to\Pi A$ for groupoid completion.
For each $a\in A$ let
\[
 C_a=\El_A\bigl(\operatorname{Hom}_{\Pi A}(a,\ell(-))\bigr).
\]
Thus an object is $(x,\gamma:a\to\ell x)$, and a morphism
$(x,\gamma)\to(y,\delta)$ is $f:x\to y$ in $A$ with
$\ell(f)\gamma=\delta$. The checked R6 theorem says, for every
fixed finite $n\geq1$, that homotopy colimits over $A$ preserve
$n$-types if and only if all these categories $C_a$ are filtered.
We re-read the supplied proof in TOTEX D.9--D.10, including its
simply connected case and the universal-cover comparison. The
covering comparison is the pullback of $BC_a\to BA$ along
$\hocolim_A X\to BA$, modeled by $\hocolim_{C_a}X|_{C_a}$:
over a bar simplex a lift is a choice of the initial path class
$\gamma$, and determines all remaining lifts. This specifies the
comparison used, rather than treating passage to a cover as an
automatic transfer principle.

In particular all the positive finite truncation classes coincide.
Their intersection with componentwise totality is componentwise
filteredness. On nonempty connected categories a separator between
totality and exactness would therefore be nonsifted and would
preserve no positive finite truncation level. These exclusions were
already established in R6; the present manuscript audit does not
claim them as new results.
\end{remark}

\begin{proposition}[Known corrections: the two types and XII's exact quantifiers]
\label{D:manuscript-types}
The following are the corrected readings of the supplied edition.
\begin{enumerate}[label=(\roman*)]
\item At II ms.\ 36, for $f:X\to Y$, the condition printed as
connectedness of $y\backslash Y$ must be connectedness of
$y\backslash X=(y\downarrow f)$, for each $y\in Y$.
\item At XII ms.\ 88 the displayed formula must be
$U'/t'\simeq U'/H'$, with $H'=p_!H$, rather than $U/H'$.
\item XII Definition~1(a), ms.\ 88--89, concerns a presheaf
$H:U^{\mathrm{op}}\to\Set$ and all functors $p:U\to U'$.
It requires the presheaf $p_!H:(U')^{\mathrm{op}}\to\Set$ to
have $W$-aspherical category of elements for every such $p$.
Definition~1(b) applies this definition to the terminal presheaf
on $U$. It is this category totality, with $W=W_\infty$, that
is compared with exactness here.
\end{enumerate}
\end{proposition}
\begin{proof}
The category $y\backslash Y$ always has initial object $\id_y$;
its connectedness cannot express finality of $f$. In contrast,
$(y\downarrow f)$ has objects $(x,\alpha:y\to f(x))$ and arrows
$r:x\to x'$ with $f(r)\alpha=\alpha'$. Nonemptiness and
connectedness of all these comma categories are precisely ordinary
finality. This is the typed form required by the question on
II ms.\ 36. Under the R1 equivalences supplied in the prompt,
totality says that every ordinary final image has weakly
contractible classifying space; it is not merely the finality of
one particular map.

For XII we keep the presheaf variance explicit. Write
$\pi_H:U/H\to U$ for the category with objects $(u,h)$,
$h\in H(u)$, and arrows $r:(u,h)\to(v,k)$ satisfying
$H(r)k=h$. The adjunction
\[
 p_!:\Set^{U^{\mathrm{op}}}
 \ \rightleftarrows\ \Set^{(U')^{\mathrm{op}}}:p^*
\]
has $p^*F'=F'\circ p^{\mathrm{op}}$ and the pointwise formula
\[
 (p_!H)(u')=\pi_0\bigl(u'\downarrow(p\circ\pi_H)\bigr).
\]
This types the edition's formula $\pi_0(u'\backslash U/H)$.
Since $H'$ is a presheaf on $U'$, only $U'/H'$ is its category
of elements. Taking $H=e_U$ gives
$(p_!e_U)(u')=\pi_0(u'\downarrow p)$, exactly Definition~1(b).
In particular there is a canonical isomorphism
$(U/H)^{\mathrm{op}}\cong\El_{U^{\mathrm{op}}}H$, with $H$
on the right treated as a covariant functor on $U^{\mathrm{op}}$.
This explicitly relates the two variance conventions.

The presheaf assertion in Definition~1(a) quantifies over
extensions along maps out of its given base $U$. It does not,
by that wording alone, quantify over all actions on $U/H$.
The supplied PRETOT presheaf counterexample prevents replacing
that restricted statement by category totality of $U/H$.
Consequently neither the corrected typography nor that
presheaf formulation supplies the missing total-to-exact
implication. The two open adjacent comparisons remain exactly
those recorded in Proposition~\ref{D:class-diagram}.
\end{proof}
